\documentclass[12pt, a4paper]{article}
\usepackage[T2A]{fontenc}
\usepackage[utf8]{inputenc}
\usepackage[russian]{babel}
\usepackage{amsmath, amsfonts, amssymb}
\usepackage{geometry}
\geometry{left=30mm, right=15mm, top=20mm, bottom=20mm}
\usepackage{graphicx}
\usepackage{listings}
\usepackage{algorithm}
\usepackage{algpseudocode}
\usepackage{tocloft}
\usepackage{indentfirst}
\usepackage{setspace}
\usepackage{booktabs}
\usepackage{tikz}
\usetikzlibrary{shapes,arrows,positioning}
\usepackage{hyperref}

\onehalfspacing

\hypersetup{
    colorlinks=true,
    linkcolor=blue,
    filecolor=magenta,      
    urlcolor=cyan,
    pdftitle={Курсовая работа: BbValidator},
    pdfpagemode=FullScreen,
}

\renewcommand{\cftsecleader}{\cftdotfill{\cftdotsep}}

% Настройка lstlisting для псевдокода
\lstset{
    basicstyle=\small\ttfamily,
    breaklines=true,
    frame=tb,
    numbers=left,
    numberstyle=\tiny,
    keywordstyle=\color{blue}\bfseries,
    commentstyle=\color{gray}\itshape,
    stringstyle=\color{green!60!black},
    language=Python,
    showstringspaces=false
}

\begin{document}

\hypersetup{pageanchor=false}
% --- ТИТУЛЬНЫЙ ЛИСТ ---
\begin{titlepage}
\begin{center}
МИНИСТЕРСТВО НАУКИ И ВЫСШЕГО ОБРАЗОВАНИЯ РОССИЙСКОЙ ФЕДЕРАЦИИ\\
\vspace{0.5cm}
ФЕДЕРАЛЬНОЕ ГОСУДАРСТВЕННОЕ БЮДЖЕТНОЕ ОБРАЗОВАТЕЛЬНОЕ УЧРЕЖДЕНИЕ\\
ВЫСШЕГО ОБРАЗОВАНИЯ\\
\vspace{1.5cm}
{\large\bf Кафедра биоинформатики и системной биологии}\\
\vspace{2.5cm}
{\Large\bf КУРСОВАЯ РАБОТА}\\
\vspace{1cm}
{\large на тему:}\\
\vspace{0.5cm}
{\Large\bf Разработка и исследование нейросетевой модели для оценки свертываемости сгенерированных белковых структур на основе гибридной графово-трансформерной архитектуры}\\
\end{center}

\vspace{3.5cm}

\begin{flushright}
\begin{minipage}{0.45\textwidth}
\begin{spacing}{1.1}
{\bf Выполнил:}\\
Студент курсовой группы\\
Иванов Иван Иванович\\
\\
{\bf Научный руководитель:}\\
д.б.н., профессор\\
Петров Петр Петрович
\end{spacing}
\end{minipage}
\end{flushright}

\vspace{\fill}

\begin{center}
Москва, 2026
\end{center}
\end{titlepage}

\newpage
\hypersetup{pageanchor=true}

% --- АННОТАЦИЯ ---
\section*{Аннотация}
Данная курсовая работа посвящена разработке и исследованию глубокой нейросетевой модели \texttt{BbValidator} (ProteinScoreModel), предназначенной для оценки свертываемости (foldability) трехмерных белковых структур (backbone), сгенерированных с помощью современных генеративных алгоритмов (таких как RFDiffusion, Chroma и др.). Модель принимает на вход пространственные координаты атомов остова ($N$, $C_\alpha$, $C$) белковой цепи и за один прямой проход решает комплексную мультизадачную проблему: предсказывает бинарный статус свертываемости, оценивает среднеквадратичное отклонение (RMSD) относительно гипотетического нативного состояния, выявляет стерические конфликты (clashes) и классифицирует доминирующий тип ошибки геометрии. Плотность водородных связей остова при этом используется как входной биофизический признак, но не как отдельная предсказательная задача. 

Архитектура модели основана на гибридном энкодере, объединяющем графовые нейронные сети (MPNN) на базе библиотеки PyTorch Geometric для анализа локальных контактов и разреженных графов взаимодействия, и полносвязный трансформер (Transformer Encoder) для учета глобального контекста и дальнодействующих взаимодействий по последовательности. Для оптимизации модели используется взвешивание потерь на основе гомоскедастической неопределенности задач (Homoscedastic Task Uncertainty), а для получения калиброванной оценки уверенности предсказаний на этапе инференса реализован подход Монте-Карло Дропаута (MC-Dropout).

\newpage

% --- СОДЕРЖАНИЕ ---
\tableofcontents
\newpage

% --- ВВЕДЕНИЕ ---
\section{Введение}

\subsection{Мотивация работы}
В последние годы методы вычислительной биологии и биоинформатики совершили качественный скачок благодаря применению глубокого машинного обучения. Модели прогнозирования пространственной структуры белков по их аминокислотным последовательностям, такие как AlphaFold и ESMFold, революционизировали структурную биологию. Вслед за этим началось активное развитие обратной задачи — \textit{дизайна белков de novo} (de novo protein design), где целью является создание новых белковых молекул с заданными функциями, топологией или свойствами связывания мишеней. 

Современные генеративные подходы, основанные на диффузионных моделях (например, RFDiffusion) и авторегрессионных генераторах (например, Chroma), способны генерировать миллионы пространственных структур белкового остова (backbone) за крайне короткое время. Однако, существенная часть сгенерированных структур оказывается физически нереализуемой: они содержат некорректную локальную геометрию, стерические столкновения атомов, петли аномальной длины или неупакованные гидрофобные ядра. 

Влажное лабораторное тестирование (in vitro) синтезированных белков является чрезвычайно дорогим и времязатратным процессом. В связи с этим возникает критическая необходимость в создании высокопроизводительных, точных и физически обоснованных вычислительных фильтров (валидаторов), способных быстро оценивать жизнеспособность сгенерированного остова до этапа дизайна боковых цепей и физического синтеза. Разработка такого валидатора, способного не просто давать бинарный ответ, но и детально интерпретировать причины ошибок и калибровать собственную уверенность предсказания, составляет основную практическую и научную ценность данной работы.

\subsection{Постановка задачи}
Целью работы является разработка и исследование нейросетевой модели оценки качества пространственной структуры остова белка. 
Входными данными для модели являются координаты атомов остова $X \in \mathbb{R}^{L \times 3 \times 3}$, где $L$ — длина белковой цепи, а размерность $3 \times 3$ соответствует пространственным координатам $(x, y, z)$ трех ключевых атомов каждого аминокислотного остатка: азота пептидной группы ($N$), альфа-углерода ($C_\alpha$) и карбонильного углерода ($C$).

Модель должна решать следующие задачи одновременно:
\begin{enumerate}
    \item \textbf{Бинарная оценка свертываемости}: оценка вероятности $P(\mathrm{Foldable}) \in [0, 1]$ стабильности пространственного остова.
    \item \textbf{Регрессия геометрического соответствия (RMSD)}: оценка среднеквадратичного отклонения предсказанной структуры от ее стабильного нативного аналога.
    \item \textbf{Оценка стерической плотности (Clashes)}: предсказание плотности стерических столкновений виртуальных атомов $C_\beta$.
    \item \textbf{Классификация типа дефекта (Failure Mode)}: распределение вероятностей по шести классам: отсутствие ошибок (\texttt{Ok}), легко обнаруживаемые дефекты (\texttt{Easy}), нарушения упаковки ядра или ложная компактность (\texttt{Core/Compact}), почти нативные искажения (\texttt{NearNative}), пограничные локальные дефекты (\texttt{Borderline}); шестой выход зарезервирован под структуры неизвестного происхождения (\texttt{Unknown}).
\end{enumerate}

Плотность водородных связей остова (число связей $O_i \ldots N_j$ с $j \ge i+3$ на остаток) вычисляется физическим фронтендом и подается в модель как входной признак, однако отдельная регрессионная голова для него не используется: предварительные эксперименты показали высокую шумовую составляющую этого таргета.

Кроме того, модель должна предоставлять количественную оценку эпистемической неопределенности (epistemic uncertainty) для калибровки предсказаний на малоизученных типах сгенерированных топологий.

\newpage

% --- ТЕОРЕТИЧЕСКИЕ ОСНОВЫ И ОПРЕДЕЛЕНИЯ ---
\section{Теоретические основы и определения}

\subsection{Пространственная геометрия остова белка}
Остов белковой цепи состоит из повторяющихся звеньев $-[N - C_\alpha - C]_i-$, соединенных ковалентными пептидными связями. 
Для математического описания конформации остова используются внутренние координаты — длины связей, валентные углы и торсионные (двугранные) углы. Поскольку длины связей и валентные углы флуктуируют около жестко фиксированных средних значений, конформация остова практически полностью детерминируется тремя двугранными углами для каждого остатка $i$: $\phi_i$ (phi), $\psi_i$ (psi) и $\omega_i$ (omega).

Определим координаты положений атомов пептидной группы остатка $i$ как векторы $\mathbf{r}_{N, i}$, $\mathbf{r}_{C_\alpha, i}$ и $\mathbf{r}_{C, i}$ в трехмерном пространстве $\mathbb{R}^3$.
Торсионный угол определяется для четырех последовательно соединенных атомов $A-B-C-D$ с координатами $\mathbf{r}_A$, $\mathbf{r}_B$, $\mathbf{r}_C$, $\mathbf{r}_D$ как угол между нормалями к полуплоскостям $ABC$ и $BCD$.
С целью исключения сингулярностей и учета направления скручивания, вводится точная векторная схема вычисления двугранного угла (метод Грама-Шмидта), реализованная в программном комплексе:
\begin{align}
    \mathbf{b}_0 &= \mathbf{r}_A - \mathbf{r}_B \\
    \mathbf{b}_1 &= \mathbf{r}_C - \mathbf{r}_B \\
    \mathbf{b}_2 &= \mathbf{r}_D - \mathbf{r}_C
\end{align}
Определим единичный направляющий вектор центральной связи как:
\begin{equation}
    \mathbf{n}_1 = \frac{\mathbf{b}_1}{\|\mathbf{b}_1\|_2}
\end{equation}
Спроецируем крайние векторы связей $\mathbf{b}_0$ и $\mathbf{b}_2$ на плоскость, ортогональную $\mathbf{n}_1$:
\begin{align}
    \mathbf{v} &= \mathbf{b}_0 - (\mathbf{b}_0 \cdot \mathbf{n}_1)\mathbf{n}_1 \\
    \mathbf{w} &= \mathbf{b}_2 - (\mathbf{b}_2 \cdot \mathbf{n}_1)\mathbf{n}_1
\end{align}
Тогда значение двугранного угла $\theta(A, B, C, D) \in [-180^\circ, 180^\circ]$ вычисляется в градусах как:
\begin{equation}
    \theta(A, B, C, D) = \operatorname{atan2}\left( (\mathbf{n}_1 \times \mathbf{v}) \cdot \mathbf{w}, \mathbf{v} \cdot \mathbf{w} \right) \cdot \frac{180^\circ}{\pi}
\end{equation}

Двугранные углы остова вычисляются по формулам:
\begin{align}
    \phi_i &= \theta(\mathbf{r}_{C, i-1}, \mathbf{r}_{N, i}, \mathbf{r}_{C_\alpha, i}, \mathbf{r}_{C, i}) \\
    \psi_i &= \theta(\mathbf{r}_{N, i}, \mathbf{r}_{C_\alpha, i}, \mathbf{r}_{C, i}, \mathbf{r}_{N, i+1}) \\
    \omega_i &= \theta(\mathbf{r}_{C_\alpha, i-1}, \mathbf{r}_{C, i-1}, \mathbf{r}_{N, i}, \mathbf{r}_{C_\alpha, i})
\end{align}

\subsection{Карта Рамачандрана и геометрические выбросы}
Вследствие стерических ограничений (исключения перекрывания ван-дер-ва\-аль\-со\-вых радиусов атомов боковых цепей и остова) углы $\phi$ и $\psi$ не могут принимать произвольные значения. Допустимые области пространства $\phi-\psi$ отображаются на карте Рамачандрана. Модель проверяет соответствие углов трем структурным зонам:
\begin{itemize}
    \item \textbf{Область $\alpha$-спиралей}: $\phi \in (-90^\circ, -30^\circ)$ и $\psi \in (-80^\circ, -10^\circ)$.
    \item \textbf{Область $\beta$-листов}: $\phi \in (-170^\circ, -50^\circ)$ и $\psi \in (80^\circ, 180^\circ)$.
    \item \textbf{Область левозакрученных $\alpha$-спиралей ($\mathrm{L}$-спирали)}: $\phi \in (30^\circ, 90^\circ)$ и $\psi \in (20^\circ, 80^\circ)$.
\end{itemize}

Остаток считается геометрическим выбросом (Ramachandran outlier), если его пара углов $(\phi_i, \psi_i)$ выпадает за пределы всех трех указанных областей:
\begin{equation}
    \text{Outlier}_i = \neg \left( \text{in\_alpha}_i \lor \text{in\_beta}_i \lor \text{in\_lalpha}_i \right) \land \text{phi\_mask}_i \land \text{psi\_mask}_i \land M_i
\end{equation}
где $M_i$ — маска валидности остатка, а $\text{phi\_mask}_i$ и $\text{psi\_mask}_i$ — маски крайних остатков ($i > 0$ и $i < N-1$), исключающие ложные выбросы на концах цепи из-за паддинга. Выбросы Рамачандрана служат надежным индикатором локального механического напряжения цепи или грубых ошибок в сгенерированной топологии.

\subsection{Восстановление виртуальных атомов и физические прокси}
Для проведения физико-химической оценки белкового остова без явного построения боковых групп, модель анализирует координаты виртуальных атомов: бета-углерода ($C_\beta$) и карбонильного кислорода ($O$).

\paragraph{Виртуальный атом $C_\beta$.}
Для всех аминокислот, кроме ахирального глицина, атом $C_\beta$ жестко ориентирован относительно пептидного остова. Его положение восстанавливается через направляющие векторы $\mathbf{v}_{N} = \mathbf{r}_{N, i} - \mathbf{r}_{C_\alpha, i}$ и $\mathbf{v}_{C} = \mathbf{r}_{C, i} - \mathbf{r}_{C_\alpha, i}$:
\begin{align}
    \mathbf{u}_{bis} &= \operatorname{normalize}\left(\operatorname{normalize}(\mathbf{v}_N) + \operatorname{normalize}(\mathbf{v}_C)\right) \\
    \mathbf{u}_{norm} &= \operatorname{normalize}\left(\operatorname{normalize}(\mathbf{v}_N) \times \operatorname{normalize}(\mathbf{v}_C)\right) \\
    \mathbf{d}_{C_\beta} &= -0.58 \cdot \mathbf{u}_{bis} + 0.81 \cdot \mathbf{u}_{norm} \\
    \mathbf{r}_{C_\beta, i} &= \mathbf{r}_{C_\alpha, i} + 1.53 \text{ \AA} \cdot \operatorname{normalize}(\mathbf{d}_{C_\beta})
\end{align}
Стерические конфликты (clashes) регистрируются, если евклидово расстояние между виртуальными атомами $C_\beta$ несет нефизический характер:
\begin{equation}
    \text{Clash}_{ij} = [\|\mathbf{r}_{C_\beta, i} - \mathbf{r}_{C_\beta, j}\|_2 < 3.5 \text{ \AA}] \quad \forall |i-j| > 0
\end{equation}

\paragraph{Виртуальный карбонильный кислород $O$.}
Положение атома кислорода необходимо для детекции водородных связей, стабилизирующих вторичную структуру. Его пространственное положение находится из планарного приближения карбонильной группы с использованием направления связи следующего остатка $\mathbf{v}_{CN} = \mathbf{r}_{N, i+1} - \mathbf{r}_{C, i}$ и направления $\mathbf{u}_{CO} = \mathbf{r}_{C_\alpha, i} - \mathbf{r}_{C, i}$:
\begin{align}
    \mathbf{d}_{O} &= \operatorname{normalize}\left(\operatorname{normalize}(\mathbf{u}_{CO}) + \operatorname{normalize}(\mathbf{v}_{CN})\right) \\
    \mathbf{r}_{O, i} &= \mathbf{r}_{C, i} + 1.23 \text{ \AA} \cdot \mathbf{d}_{O}
\end{align}
Водородная связь фиксируется при выполнении геометрического условия сближения донорно-акцепторных атомов:
\begin{equation}
    \text{HBond}_{ij} = [\|\mathbf{r}_{O, i} - \mathbf{r}_{N, j}\|_2 < 3.5 \text{ \AA}] \quad \forall j - i \ge 3
\end{equation}

\subsection{Биофизические дескрипторы свертываемости}
Помимо точечных геометрических углов, модель использует макроскопические признаки упаковки:
\begin{enumerate}
    \item \textbf{Плотность упаковки ядра (Contact Density)}: число удаленных по последовательности остатков в радиальных окрестностях $6, 8, 10\,\text{\AA}$:
    \begin{equation}
        D_{r, i} = \sum_{j: |i-j| > 3} [\|\mathbf{r}_{C_\alpha, i} - \mathbf{r}_{C_\alpha, j}\|_2 < r] \cdot M_j
    \end{equation}
    \item \textbf{Относительная экспонированность (Relative Burial)}: положение остатка относительно центроида глобулы $\mathbf{r}_{center}$, нормированное на теоретический радиус $R_{expected} = 2.2 \cdot N_{eff}^{1/3}$:
    \begin{equation}
        B_i = \frac{\|\mathbf{r}_{C_\alpha, i} - \mathbf{r}_{center}\|_2}{R_{expected}}
    \end{equation}
    \item \textbf{Локальный изгиб цепи (Local Bending)}: прокси-показатель кривизны петель, вычисляемый как смещение между остатками с шагом 4:
    \begin{equation}
        K_i = \|\mathbf{r}_{C_\alpha, i+2} - \mathbf{r}_{C_\alpha, i-2}\|_2
    \end{equation}
    \item \textbf{Фрагментный признак (PDB-Fragment Similarity)}: 16-мерный дескриптор, получаемый линейным проецированием вектора попарных расстояний внутри локального окна размером $K=9$ остатков.
    
    Для каждого остатка $i \in \{1, \dots, N\}$ выделяется скользящее окно из 9 последовательных остатков, центрированное относительно $i$. Для обработки краев последовательности применяется дублирующий паддинг (replicate padding) на 4 элемента слева и справа. Пусть координаты атомов $C_\alpha$ в окне для остатка $i$ равны $\{\mathbf{x}_{i, 1}, \mathbf{x}_{i, 2}, \dots, \mathbf{x}_{i, 9}\}$, где $\mathbf{x}_{i, a} \in \mathbb{R}^3$.
    
    Для каждого такого окна рассчитывается матрица попарных расстояний $\mathbf{D}^{(i)} \in \mathbb{R}^{9 \times 9}$:
    \begin{equation}
        D^{(i)}_{ab} = \|\mathbf{x}_{i, a} - \mathbf{x}_{i, b}\|_2, \quad a, b \in \{1, \dots, 9\}
    \end{equation}
    Из полученной симметричной матрицы извлекаются элементы строго над главной диагональю, формирующие 36-мерный вектор попарных расстояний остатков в окне $\mathbf{d}_{pairs, i} \in \mathbb{R}^{36}$:
    \begin{equation}
        \mathbf{d}_{pairs, i} = \left[ D^{(i)}_{ab} \right]_{1 \le a < b \le 9}
    \end{equation}
    Для сжатия признаков и снижения шума вектор $\mathbf{d}_{pairs, i}$ линейно проецируется в 16-мерное пространство признаков $\mathbf{P}_{frag, i} \in \mathbb{R}^{16}$ с помощью PCA-матрицы проекции $\mathbf{W}_{PCA} \in \mathbb{R}^{16 \times 36}$ (что соответствует слою \texttt{Linear(36, 16, bias=False)} в классе \texttt{FoldabilityProxies} проекта):
    \begin{equation}
        \mathbf{P}_{frag, i} = \mathbf{W}_{PCA} \cdot \mathbf{d}_{pairs, i}
    \end{equation}
    Матрица $\mathbf{W}_{PCA}$ рассчитывается заранее методом главных компонент (PCA) на основе наборов структур реальных белков из базы PDB и фиксируется (замораживается) в процессе обучения основной модели.
    
    \textit{Физический смысл:} 36 расстояний полностью определяют жесткую геометрию 9-остаточного фрагмента (его пространственную укладку). Линейная проекция сжимает эту информацию до 16 главных признаков, кодирующих, насколько локальная форма похожа на стандартные природные элементы вторичной структуры ($\alpha$-спирали, $\beta$-листы, петельные переходы).
    
    \paragraph{Реализация в проекте.}
    Вычисление локальной фрагментной схожести программно реализовано в модуле препроцессинга \href{file:///home/pc/PycharmProjects/BbValidator/preprocess/foldability_features.py}{\texttt{preprocess/foldability\_features.py}} в рамках класса \texttt{FoldabilityProxies}. Извлечение попарных расстояний внутри скользящих окон выполняется оптимизированным векторным методом с использованием операций библиотеки \texttt{PyTorch} без явных циклов: дополнение краев последовательности через \texttt{torch.nn.functional.pad}, формирование окон через \texttt{unfold} и вычисление попарных расстояний через \texttt{torch.cdist}.
    
    В Листинге~\ref{lst:pca_frag} представлен соответствующий фрагмент кода реализации методов \texttt{\_get\_local\_pairwise\_distances} и \texttt{forward} класса \texttt{FoldabilityProxies}.

\begin{lstlisting}[language=Python, caption={Реализация вычисления PDB-Fragment Similarity (из \texttt{foldability\_features.py})}, label={lst:pca_frag}]
# FoldabilityProxies class in preprocess/foldability_features.py
class FoldabilityProxies(nn.Module):
    def __init__(self, contact_threshold=8.0, seq_sep=3, pca_components=16, fragment_size=9):
        super().__init__()
        self.fragment_size = fragment_size
        self.frag_pairs = (fragment_size * (fragment_size - 1)) // 2
        # Upper triangular indices to extract unique pairs
        triu = torch.triu_indices(fragment_size, fragment_size, offset=1)
        self.register_buffer("triu_indices", triu)
        # Linear layer for PCA projection (weights are frozen)
        self.pca_proj = nn.Linear(self.frag_pairs, pca_components, bias=False)

    def _get_local_pairwise_distances(self, coords: torch.Tensor) -> torch.Tensor:
        B, N, _ = coords.shape
        K = self.fragment_size
        pad_left = K // 2
        pad_right = K - 1 - pad_left
        
        # Replicate padding at the sequence boundaries along residue axis
        coords_padded = F.pad(
            coords.transpose(1, 2), (pad_left, pad_right), mode="replicate"
        ).transpose(1, 2)
        
        # Extract sliding windows of size K along residue axis (dimension=1)
        windows = coords_padded.unfold(dimension=1, size=K, step=1)  # [B, N, 3, K]
        windows = windows.permute(0, 1, 3, 2).contiguous()  # [B, N, K, 3]
        
        # Compute pairwise Euclidean distances within each window
        dist_mat = torch.cdist(windows, windows)  # [B, N, K, K]
        
        # Extract upper triangular part (36 pairwise distances)
        flat_distances = dist_mat[
            :, :, self.triu_indices[0], self.triu_indices[1]
        ]  # [B, N, frag_pairs]
        return flat_distances

    def forward(self, ca_coords: torch.Tensor, mask=None, dist_mat=None):
        # ... (compute Contact Density and Relative Burial)
        # 3. PDB-Fragment Similarity
        fragment_features = self._get_local_pairwise_distances(ca_coords) # [B, N, 36]
        pca_projection = self.pca_proj(fragment_features)                # [B, N, 16]
        # ...
        return {
            # ...
            "pca_projection": pca_projection * mask.float().unsqueeze(-1),
        }
\end{lstlisting}
\end{enumerate}

\subsection{Попарные признаки и построение графа контактов}
Белок представляет собой связную физическую структуру, целостность которой определяется не только локальной геометрией отдельных остатков, но и дальнодействующими взаимодействиями в пространстве (такими как водородные связи, гидрофобные контакты и электростатические силы). Чтобы закодировать эти трехмерные пространственные взаимосвязи, модель вычисляет набор попарных признаков взаимодействия остатков $(i, j)$ и строит разреженный контактный граф.

Использование полных попарных матриц взаимодействия размера $N \times N$ приводит к квадратичной вычислительной сложности $O(N^2)$, что делает обработку длинных белков крайне неэффективной и требовательной к ресурсам GPU. Для оптимизации вычислений строится разреженный ориентированный граф ближайших соседей (kNN-граф, $k=16$). Для каждого остатка определяются ребра к $16$ пространственно ближайшим соседям по расстоянию между атомами $C_\alpha$. Это позволяет снизить сложность локальной графовой свертки до линейной $O(N)$.

Помимо узловых признаков, для каждой пары остатков $(i, j)$ рассчитывается вектор парных признаков $\mathbf{x}_{pair, ij} \in \mathbb{R}^{20}$, состоящий из следующих компонентов:
\begin{itemize}
    \item \textbf{RBF-разложение расстояния (Radial Basis Function Expansion)}:
    Расстояние $d_{ij} = \|\mathbf{r}_{C_\alpha, i} - \mathbf{r}_{C_\alpha, j}\|_2$ между альфа-углеродами остатков $i$ и $j$ кодируется вектором из $B=16$ значений с помощью радиально-базисных (гауссовых) функций, равномерно распределенных на интервале от $d_{min} = 2.0\,\text{\AA}$ до $d_{max} = 22.0\,\text{\AA}$:
    \begin{equation}
        \mathrm{RBF}_{ij, k} = \exp\left(-\gamma (d_{ij} - \mu_k)^2\right), \quad k \in \{1, \dots, B\}
    \end{equation}
    где центры функций $\mu_k$ задаются линейной сеткой:
    \begin{equation}
        \mu_k = d_{min} + (k - 1) \cdot \frac{d_{max} - d_{min}}{B - 1}
    \end{equation}
    а коэффициент ширины ядра $\gamma$ определяется шагом сетки $\Delta = \frac{d_{max} - d_{min}}{B}$:
    \begin{equation}
        \gamma = \frac{1}{\Delta^2} = \frac{1}{(1.25\,\text{\AA})^2} = 0.64\,\text{\AA}^{-2}
    \end{equation}
    
    \textit{Справка по RBF-разложению:} В задачах трехмерного анализа молекулярных структур прямое использование скалярного расстояния $d_{ij}$ в качестве входного признака малоэффективно, так как нейронные сети плохо аппроксимируют высокочастотные локальные изменения расстояний с помощью одного вещественного числа. RBF-разложение (или «плавное бинирование», soft binning) представляет собой отображение расстояния в высокомерный вектор признаков. Вместо жесткого разбиения на интервалы (гистограммы), где переход через границу ячейки разрывен, RBF обеспечивает гладкое, дифференцируемое представление: при изменении расстояния плавно меняются активации сразу нескольких соседних гауссовых каналов. Это стабилизирует градиенты при обучении и позволяет MPNN-слоям эффективнее улавливать пространственные контакты на разных масштабах (от ван-дер-ваальсовых взаимодействий до дальнодействующих петель).
    
    \item \textbf{Бинарный контактный флаг}: $[d_{ij} < 8.0\,\text{\AA}]$, отражающий непосредственный физический контакт между остатками.
    \item \textbf{Признаки взаимного положения по последовательности}:
    \begin{itemize}
        \item Нормализованное расстояние по цепи: $\min(1.0, |i-j| / 32.0)$.
        \item Индикатор непосредственного соседства по цепи: $[|i-j| == 1]$.
        \item Индикатор локального окна: $[|i-j| \le 4]$.
    \end{itemize}
\end{itemize}

\paragraph{Реализация в проекте.}
Процесс формирования попарных признаков и фильтрации kNN-графа реализован в модуле \href{file:///home/pc/PycharmProjects/BbValidator/preprocess/pair_features.py}{\texttt{preprocess/pair\_features.py}} в рамках классов \texttt{RBFExpansion} и \texttt{PairFeatureBuilder}.

\subsection{Гомоскедастическая неопределенность задач}
При одновременном обучении модели на четыре различные задачи (бинарный статус свертываемости, регрессия RMSD в логарифмической шкале $\ln(1+\text{RMSD})$, плотность стерических конфликтов и тип ошибки) традиционное суммирование потерь $\mathcal{L} = \sum w_i \mathcal{L}_i$ неустойчиво. Для автоматической балансировки весов лоссов используется байесовский подход Kendall et al. \cite{kendall2018}, основанный на гомоскедастической неопределенности каждой задачи:
\begin{equation}
    \mathcal{L}(W, s_1, \dots, s_T) = \sum_{t=1}^T \left( \frac{1}{2} e^{-s_t} \mathcal{L}_t(W) + \frac{1}{2} s_t \right)
\end{equation}
где $s_t = \ln \sigma_t^2$ — вещественные обучаемые параметры модели, регулирующие веса потерь $\tau_t = e^{-s_t}$. Параметр $s_t$ штрафует модель за завышение дисперсии (неопределенности), предотвращая стремление лосса к минус бесконечности.

\subsection{Эпистемическая неопределенность и Монте-Карло Дропаута}
Для предотвращения пропусков аномальных структур, далеких от обучающего распределения, в фазе инференса применяется метод Монте-Карло Дропаута \cite{gal2016}. Слои Dropout принудительно переводятся в режим обучения (\texttt{train()}), в то время как слои LayerNorm остаются замороженными (\texttt{eval()}).
Модель осуществляет $M$ прямых проходов для одного и того же графа. Математические ожидания предсказаний вероятности фолдинга $\bar{p}$ и меры неопределенности $U$ (выражающей эпистемическую неуверенность модели) рассчитываются следующим образом:
\begin{align}
    \bar{p} &= \frac{1}{M} \sum_{m=1}^M \sigma(l_m) \\
    U &= \operatorname{Var}\left(\{ \sigma(l_m) \}_{m=1}^M \right) = \frac{1}{M} \sum_{m=1}^M \left( \sigma(l_m) - \bar{p} \right)^2
\end{align}
где $l_m$ — логит фолдинга на проходе $m$, $\sigma(z) = 1/(1+e^{-z})$. При $U > U_{thr}$ структура помечается как сомнительная (высокая дисперсия предсказания).


\newpage

% --- АРХИТЕКТУРА И СТРУКТУРА ПРОЕКТА ---
\section{Архитектура и структура проекта}

\subsection{Модульная структура репозитория}
Проект реализован на языке Python с использованием фреймворка PyTorch. Архитектура разбита на логические компоненты, представленные в таблице \ref{tab:structure}.

\begin{table}[h]
\centering
\caption{Структура исходного кода проекта}
\label{tab:structure}
\begin{tabular}{p{5.5cm}p{9.5cm}}
\toprule
\textbf{Путь к модулю} & \textbf{Функциональное назначение} \\
\midrule
\raggedright\texttt{dataset/\allowbreak dataloader.py} & Чтение HDF5, SE(3)-аугментация, динамический батчинг. \\
\raggedright\texttt{preprocess/\allowbreak biophys\_\allowbreak frontend.py} & Вычисление физических признаков остова без градиентов. \\
\raggedright\texttt{preprocess/\allowbreak geometry\_\allowbreak features.py} & Расчет двугранных углов, виртуальных атомов, клэшей и водородных связей. \\
\raggedright\texttt{preprocess/\allowbreak foldability\_\allowbreak features.py} & Плотность упаковки, экспонированность, фрагментные PCA признаки. \\
\raggedright\texttt{preprocess/\allowbreak pair\_\allowbreak features.py} & Построение признаков пар и kNN-графа. \\
\raggedright\texttt{model/\allowbreak encoder.py} & Гибридный кодировщик: слои MPNN + Transformer Encoder. \\
\raggedright\texttt{model/\allowbreak heads\_\allowbreak loss.py} & Головы задач, пулинг внимания, динамический лосс, MC-Dropout. \\
\raggedright\texttt{train\_model.py} & Главный цикл обучения модели с bfloat16 AMP и GradScaler. \\
\raggedright\texttt{inference.py} & Скрипт инференса для оценки PDB-структур. \\
\bottomrule
\end{tabular}
\end{table}

\subsection{Конвейер обработки данных и архитектура сети}
Модель обрабатывает информацию согласно следующему конвейеру (см. рис. \ref{fig:architecture}):
\begin{figure}[h]
\centering
\begin{tikzpicture}[node distance=1.5cm, auto]
    \tikzstyle{block} = [rectangle, draw, fill=blue!10, text width=6cm, text centered, rounded corners, minimum height=1cm]
    \tikzstyle{line} = [draw, -latex', thick]
    
    \node [block] (input) {Входные координаты остова\\ $[B, L, 3, 3]$};
    \node [block, below=of input] (frontend) {BiophysicalFrontend (без градиентов)\\ Node Feats: $[B, L, 31]$ \\ Pair Feats: $[B, L, L, 20]$ + kNN};
    \node [block, below=of frontend] (mpnn) {PyG Graph Message Layer $\times 2$\\ Локальная геометрия и контакты};
    \node [block, below=of mpnn] (trans) {Transformer Encoder $\times 4$\\ Глобальный контекст по последовательности};
    \node [block, below=of trans] (pool) {Multi-Head Attention Pooling\\ Получение вектора структуры $[B, d_{model}]$};
    \node [block, below=of pool] (heads) {ProteinMultiTaskHeads\\ Выходные предсказания (4 задачи)};

    \path [line] (input) -- (frontend);
    \path [line] (frontend) -- (mpnn);
    \path [line] (mpnn) -- (trans);
    \path [line] (trans) -- (pool);
    \path [line] (pool) -- (heads);
\end{tikzpicture}
\caption{Блок-схема конвейера обработки данных модели BbValidator}
\label{fig:architecture}
\end{figure}

На этапе фронтенда физические и геометрические признаки вычисляются с использованием эффективных векторных операций в \texttt{PyTorch} и кешируются. Архитектура сети построена на гибридном принципе и включает следующие ключевые блоки:

\subsubsection{Проекционные слои (Embeddings)}
Начальные признаки узлов $X_{node} \in \mathbb{R}^{B \times N \times 31}$ и признаки пар $X_{pair} \in \mathbb{R}^{B \times N \times N \times 20}$ переводятся в размерность скрытого пространства модели $d_{model} = 192$ с помощью полносвязных проекций:
\begin{align}
    \mathbf{h}_{i}^{(0)} &= \text{LayerNorm}(\mathbf{W}_{node} \mathbf{x}_{node, i} + \mathbf{b}_{node}) \\
    \mathbf{e}_{ij}^{(0)} &= \mathbf{W}_{pair, 2} \cdot \text{GELU}(\text{LayerNorm}(\mathbf{W}_{pair, 1} \mathbf{x}_{pair, ij} + \mathbf{b}_{pair, 1})) + \mathbf{b}_{pair, 2}
\end{align}

\subsubsection{Динамическая упаковка (Dynamic Packing)}
Для повышения производительности и экономии памяти при обработке графов переменного размера, в MPNN-слоях используется упаковка тензоров. Паддинг-токены удаляются, а валидные остатки всех белков батча объединяются в единый двумерный тензор $\mathbf{H}_{packed} \in \mathbb{R}^{N_{valid} \times d_{model}}$, где $N_{valid} = \sum_{b=1}^{B} \sum_{i=1}^{N} M_{b, i}$. 
Индексы ребер в kNN-графе ($k=16$ ближайших соседей по расстоянию $C_\alpha - C_\alpha$) смещаются на накопительный офсет:
\begin{equation}
    E_{idx, global} = \text{Concat}\left(\{ E_{idx, b} + \sum_{l=1}^{b-1} N_{valid, l} \}_{b=1}^B \right)
\end{equation}
После прохождения MPNN-слоев признаки распаковываются обратно в тензор батча с заполнением позиций паддинга нулями с помощью маски.

\subsubsection{Графовые слои MPNN с GRU-обновлением}
Слой \texttt{PyGGraphMessageLayer} осуществляет локальный обмен сообщениями между соседними остатками по формулам:
\begin{equation}
    \mathbf{m}_{ij} = \mathbf{W}_{msg\_out} \cdot \text{GELU}(\mathbf{W}_{n} \mathbf{h}_i + \mathbf{W}_{n} \mathbf{h}_j + \mathbf{W}_{e} \mathbf{e}_{ij})
\end{equation}
где $\mathbf{W}_{n}$ — проекции узлов-участников ребра, а $\mathbf{W}_{e}$ — проекция признаков ребра. Сообщения агрегируются усреднением и сглаживаются регуляризационным зашумлением:
\begin{equation}
    \mathbf{a}_i = \text{Dropout}\left(\frac{1}{|N(i)|} \sum_{j \in N(i)} \mathbf{m}_{ij}\right)
\end{equation}
Для обновления скрытого состояния узла применяется GRU-подобный гейтинг (gated update):
\begin{align}
    \mathbf{z}_i &= \sigma(\mathbf{W}_z [\mathbf{h}_i; \mathbf{a}_i] + \mathbf{b}_z) \\
    \mathbf{r}_i &= \sigma(\mathbf{W}_r [\mathbf{h}_i; \mathbf{a}_i] + \mathbf{b}_r) \\
    \hat{\mathbf{h}}_i &= \tanh(\mathbf{W}_u [(\mathbf{r}_i \odot \mathbf{h}_i); \mathbf{a}_i] + \mathbf{b}_u) \\
    \mathbf{h}_i' &= (1 - \mathbf{z}_i) \odot \mathbf{h}_i + \mathbf{z}_i \odot \hat{\mathbf{h}}_i \\
    \mathbf{h}_i^{out} &= \text{LayerNorm}(\mathbf{h}_i')
\end{align}
где $\sigma(\cdot)$ — логистическая функция, $\odot$ — поэлементное произведение Адамара.

\subsubsection{Полносвязный трансформер (Transformer Encoder)}
После распаковки в исходный формат батча к тензору применяется классический Transformer Encoder ($4$ слоя с $8$ головами внимания) с механизмом Pre-LayerNorm. Механизм самовнимания (Self-Attention) позволяет учесть дальнодействующие по последовательности взаимодействия, например, укладку удаленных петель или формирование общего бета-листа. Для исключения влияния паддингов используется маска \texttt{src\_key\_padding\_mask} = $\neg M$.

\subsubsection{Агрегация признаков (Multi-Head Attention Pooling)}
Для получения единого вектора белковой структуры $\mathbf{z} \in \mathbb{R}^{d_{model}}$ применяется многоголовый пулинг с вниманием (количество голов $K=4$):
\begin{equation}
    S_{i, k} = \mathbf{w}_k \cdot \tanh(\mathbf{W}_{s} \mathbf{h}_i + \mathbf{b}_s) + b_k
\end{equation}
где $\mathbf{h}_i$ — выходной эмбеддинг остатка $i$. Веса внимания нормируются с помощью softmax по длине последовательности с маскированием паддингов (значения $S_{i,k}$ для маскированных позиций принудительно заменяются на $-\infty$ перед экспонированием):
\begin{equation}
    \alpha_{i, k} = \frac{\exp(S_{i, k})}{\sum_{j=1}^N \exp(S_{j, k})}
\end{equation}
Вектора голов вычисляются и объединяются в один плоский вектор:
\begin{align}
    \mathbf{p}_k &= \sum_{i=1}^N \alpha_{i, k} \mathbf{h}_i \in \mathbb{R}^{d_{model}} \\
    \mathbf{p}_{flat} &= [\mathbf{p}_1; \mathbf{p}_2; \dots; \mathbf{p}_K] \in \mathbb{R}^{K \cdot d_{model}}
\end{align}
Финальный вектор белка проецируется в исходную размерность:
\begin{equation}
    \mathbf{z}_{protein} = \mathbf{W}_{out} \mathbf{p}_{flat} + \mathbf{b}_{out}
\end{equation}
Этот вектор передается в головы мультизадачного вывода для финального предсказания.


\newpage

% --- АЛГОРИТМЫ И ПСЕВДОКОД ---
\section{Алгоритмы и псевдокод}

В данном разделе приведены алгоритмы ключевых модулей системы, записанные в виде структурированного псевдокода на основе оригинального программного кода проекта.

\subsection{Алгоритм экстракции геометрических признаков}
Алгоритм вычисляет локальные геометрические свойства остова: торсионные углы, выбросы Рамачандрана, плотность клэшей и количество водородных связей. Соответствует коду класса \texttt{BackboneGeometryExtractor}.

\begin{algorithm}[h]
\caption{Извлечение геометрии остова белка}
\label{alg:geometry}
\begin{algorithmic}[1]
\Statex \textbf{Вход:} Координаты остова $C \in \mathbb{R}^{B \times N \times 3 \times 3}$, Маска валидности $M \in \mathbb{B}^{B \times N}$
\Statex \textbf{Выход:} Набор локальных геометрических признаков
\State Разделить координаты на атомы: $N_{at} \gets C[:, :, 0]$, $C_\alpha \gets C[:, :, 1]$, $C_{carb} \gets C[:, :, 2]$
\State Вычислить сырые торсионные углы $\phi_{raw}, \psi_{raw}, \omega_{raw}$ по формуле двугранного угла (метод Грама-Шмидта):
\State \quad $\phi_{raw} \gets \operatorname{dihedral}(C_{carb}[:, :-1], N_{at}[:, 1:], C_\alpha[:, 1:], C_{carb}[:, 1:])$
\State \quad $\psi_{raw} \gets \operatorname{dihedral}(N_{at}[:, :-1], C_\alpha[:, :-1], C_{carb}[:, :-1], N_{at}[:, 1:])$
\State \quad $\omega_{raw} \gets \operatorname{dihedral}(C_\alpha[:, :-1], C_{carb}[:, :-1], N_{at}[:, 1:], C_\alpha[:, 1:])$
\State Выполнить паддинг и маскирование для выравнивания по размерности $N$:
\State \quad $\phi \gets \operatorname{pad\_left}(\phi_{raw}, 1)$, $\psi \gets \operatorname{pad\_right}(\psi_{raw}, 1)$, $\omega \gets \operatorname{pad\_left}(\omega_{raw}, 1)$
\State \quad $\text{phi\_mask}_i \gets [i > 0]$, $\text{psi\_mask}_i \gets [i < N-1]$ \Comment{маски исключения крайних остатков}
\State Вычислить признаки выбросов Рамачандрана:
\State \quad $\text{in\_alpha}_i \gets (\phi_i \in (-90^\circ, -30^\circ)) \land (\psi_i \in (-80^\circ, -10^\circ))$
\State \quad $\text{in\_beta}_i \gets (\phi_i \in (-170^\circ, -50^\circ)) \land (\psi_i \in (80^\circ, 180^\circ))$
\State \quad $\text{in\_lalpha}_i \gets (\phi_i \in (30^\circ, 90^\circ)) \land (\psi_i \in (20^\circ, 80^\circ))$
\State \quad $\text{outliers}_i \gets \neg(\text{in\_alpha}_i \lor \text{in\_beta}_i \lor \text{in\_lalpha}_i) \land \text{phi\_mask}_i \land \text{psi\_mask}_i \land M_i$
\State Восстановить координаты виртуального $C_\beta$:
\State \quad $\mathbf{v}_N \gets N_{at} - C_\alpha$, $\mathbf{v}_C \gets C_{carb} - C_\alpha$
\State \quad $\mathbf{u}_{bis} \gets \operatorname{normalize}\left(\operatorname{normalize}(\mathbf{v}_N) + \operatorname{normalize}(\mathbf{v}_C)\right)$
\State \quad $\mathbf{u}_{norm} \gets \operatorname{normalize}\left(\operatorname{normalize}(\mathbf{v}_N) \times \operatorname{normalize}(\mathbf{v}_C)\right)$
\State \quad $\mathbf{d}_{C_\beta} \gets -0.58 \cdot \mathbf{u}_{bis} + 0.81 \cdot \mathbf{u}_{norm}$
\State \quad $C_\beta \gets C_\alpha + 1.53\,\text{\AA} \cdot \operatorname{normalize}(\mathbf{d}_{C_\beta})$
\State Восстановить координаты виртуального кислорода $O$:
\State \quad $\mathbf{u}_{CO} \gets C_\alpha - C_{carb}$, $\mathbf{v}_{CN} \gets \operatorname{roll}(N_{at}, \text{shift}=-1) - C_{carb}$
\State \quad $\mathbf{d}_O \gets \operatorname{normalize}\left(\operatorname{normalize}(\mathbf{u}_{CO}) + \operatorname{normalize}(\mathbf{v}_{CN})\right)$
\State \quad $O \gets C_{carb} + 1.23\,\text{\AA} \cdot \mathbf{d}_O$
\State Вычислить водородные связи:
\State \quad $D_{ON} \gets \operatorname{cdist}(O, N_{at})$
\State \quad $\text{hb\_mask}_{ij} \gets (D_{ON, ij} < 3.5\,\text{\AA}) \land (j - i \ge 3) \land M_i \land M_j$
\State \quad $\text{hbond\_count}_i \gets \sum_{j} \text{hb\_mask}_{ij}$
\State Вычислить стерические клэши:
\State \quad $D_{C_\beta} \gets \operatorname{cdist}(C_\beta, C_\beta)$
\State \quad $\text{clash\_mask}_{ij} \gets (D_{C_\beta, ij} < 3.5\,\text{\AA}) \land (i \neq j) \land M_i \land M_j$
\State \quad $\text{clash\_count}_i \gets \sum_{j} \text{clash\_mask}_{ij}$
\State \Return $\{\phi, \psi, \omega, \text{outliers}, \text{hbond\_count}, \text{clash\_count}\}$
\end{algorithmic}
\end{algorithm}

\newpage

\subsection{Алгоритм экстракции признаков свертываемости}
Алгоритм извлекает признаки плотности упаковки ядра, экспонированности растворителю и фрагментной схожести с использованием PCA-проекции. Соответствует классу \texttt{FoldabilityProxies}.

\begin{algorithm}[h]
\caption{Извлечение признаков свертываемости}
\label{alg:foldability}
\begin{algorithmic}[1]
\Statex \textbf{Вход:} Координаты $C_\alpha \in \mathbb{R}^{B \times N \times 3}$, Маска $M \in \mathbb{B}^{B \times N}$, Матрица расстояний $D \in \mathbb{R}^{B \times N \times N}$
\Statex \textbf{Выход:} Прокси-признаки свертываемости
\State Инициализировать маску дальних контактов: $V_{ij} \gets (|i - j| > 3) \land M_i \land M_j$
\State Вычислить плотность контактов:
\For{каждого радиуса $r \in \{6, 8, 10\}$}
    \State $\text{density}_{r, i} \gets \sum_{j=1}^N [D_{ij} < r] \cdot V_{ij}$
\EndFor
\State Оценить относительную экспонированность:
\State \quad $\mathbf{r}_{center} \gets \left(\sum_j C_{\alpha, j} \cdot M_j\right) / \left(\sum_j M_j + 10^{-8}\right)$
\State \quad $d_{center, i} \gets \|C_{\alpha, i} - \mathbf{r}_{center}\|_2$
\State \quad $R_{expected} \gets 2.2 \cdot \left(\sum_j M_j\right)^{1/3} + 10^{-5}$
\State \quad $\text{burial}_i \gets d_{center, i} / R_{expected}$
\State Вычислить фрагментную схожесть:
\State \quad Выделить скользящие окна размера $K=9$ по $C_\alpha$ с replicate padding $(4, 4)$
\State \quad Для каждого окна собрать вектор расстояний $\mathbf{d}_{pairs} \in \mathbb{R}^{36}$ (верхний треугольник)
\State \quad Спроецировать: $P_{frag} \gets \mathbf{W}_{PCA} \cdot \mathbf{d}_{pairs}$ (где $\mathbf{W}_{PCA} \in \mathbb{R}^{16 \times 36}$)
\State Вычислить локальную кривизну (local bending):
\State \quad $K_i \gets \|C_{\alpha, i+2} - C_{\alpha, i-2}\|_2$ (с дополнением нулями для краев)
\State \Return $\{\text{density}_{6,8,10}, \text{burial}, P_{frag}, K\}$
\end{algorithmic}
\end{algorithm}

\subsection{Алгоритм гибридного кодировщика (HybridProteinEncoder)}
Этот алгоритм демонстрирует процесс интеграции графовых слоев с трансформером. Для экономии памяти используется метод градиентного чекпоинтинга на этапе графовой свертки.

\begin{algorithm}[h]
\caption{Прямой проход гибридного кодировщика}
\label{alg:encoder}
\begin{algorithmic}[1]
\Statex \textbf{Вход:} Признаки узлов $X_{node} \in \mathbb{R}^{B \times N \times 31}$, Признаки пар $X_{pair} \in \mathbb{R}^{B \times N \times N \times 20}$, Маска $M$, Рёбра графа $E_{idx}$, Атрибуты ребер $E_{attr}$
\Statex \textbf{Выход:} Закодированные эмбеддинги узлов $H \in \mathbb{R}^{B \times N \times d_{model}}$
\State Спроецировать признаки узлов: $H_{node} \gets \operatorname{LayerNorm}(\operatorname{Linear}(X_{node}))$
\State Спроецировать признаки ребер: $H_{edge} \gets \operatorname{Linear}(\operatorname{GELU}(\operatorname{LayerNorm}(\operatorname{Linear}(E_{attr}))))$
\State \textbf{Упаковка графа для MPNN (удаление паддинга):}
\State \quad $H_{packed} \gets \operatorname{Concat}(\{ H_{node}[b][M[b]] \}_{b=0}^{B-1})$
\State \quad Сдвинуть индексы ребер $E_{idx}$ каждого батча на накопительный offset
\State \quad $E_{idx, global} \gets \operatorname{Concat}(E_{idx})$
\State \quad $E_{attr, global} \gets \operatorname{Concat}(H_{edge})$
\State \textbf{Вычисление графовых сообщений (MPNN):}
\For{each layer $L_{graph} \in \text{GraphMessageLayers}$}
    \State Применить градиентный чекпоинтинг для экономии памяти GPU:
    \State $H_{packed} \gets \operatorname{checkpoint}(L_{graph}, H_{packed}, E_{idx, global}, E_{attr, global})$
\EndFor
\State \textbf{Распаковка графа обратно в тензор батча:}
\State \quad Инициализировать $H_{batch} \in \mathbb{R}^{B \times N \times d_{model}}$ нулями
\State \quad Распределить $H_{packed}$ по валидным маскированным позициям: $H_{batch}[M] \gets H_{packed}$
\State \textbf{Глобальное кодирование (Transformer):}
\State \quad $Pad_{mask} \gets \neg M$
\State \quad $H \gets \operatorname{TransformerEncoder}(H_{batch}, \text{src\_key\_padding\_mask}=Pad_{mask})$
\State \Return $H$
\end{algorithmic}
\end{algorithm}

\newpage

\subsection{Алгоритм обучения и расчета мультизадачной функции потерь}
Алгоритм вычисляет общую функцию потерь с динамическими весами на основе гомоскедастической неопределенности.

\begin{algorithm}[h]
\caption{Расчет функции потерь DynamicMultiTaskLoss}
\label{alg:loss}
\begin{algorithmic}[1]
\Statex \textbf{Вход:} Предсказания модели $P$, Таргеты батча $Y$, Параметры неопределенности $S \in \mathbb{R}^4$
\Statex \textbf{Выход:} Общий лосс $\mathcal{L}_{total}$, Словарь метрик
\State Вычислить $\mathcal{L}_{fold} \gets \operatorname{BCEWithLogits}(P_{fold\_logit}, Y_{label})$
\State Вычислить $\mathcal{L}_{rmsd} \gets \operatorname{MSE}(P_{rmsd}, \ln(1 + Y_{rmsd\_target}))$
\State Вычислить $\mathcal{L}_{steric} \gets \operatorname{MSE}(P_{steric}, Y_{steric\_target})$
\State Вычислить $\mathcal{L}_{fail} \gets \operatorname{CrossEntropy}(P_{failure\_mode}, Y_{failure\_mode\_label}, \text{weight}=W_{classes})$
\State Сгруппировать потери: $\vec{\mathcal{L}} \gets [\mathcal{L}_{fold}, \mathcal{L}_{rmsd}, \mathcal{L}_{steric}, \mathcal{L}_{fail}]$
\State Рассчитать точность (precision) для каждой задачи: $\vec{\tau} \gets \exp(-S)$
\State Рассчитать взвешенный лосс по формуле Кендалла-Гала:
\State \quad $\mathcal{L}_{total} \gets \sum_{t=1}^4 \left( 0.5 \cdot \tau_t \cdot \mathcal{L}_t + 0.5 \cdot S_t \right)$
\State \Return $\mathcal{L}_{total}, \{ \mathcal{L}_{fold}, \mathcal{L}_{rmsd}, \mathcal{L}_{steric}, \mathcal{L}_{fail}, \text{weights}=\vec{\tau} \}$
\end{algorithmic}
\end{algorithm}

\subsection{Алгоритм MC-Dropout инференса}
Алгоритм оценивает эпистемическую неопределенность структуры путем многократного прохождения через граф-трансформер с активным дропаутом.

\begin{algorithm}[h]
\caption{Оценка неопределенности методом MC-Dropout}
\label{alg:mcdropout}
\begin{algorithmic}[1]
\Statex \textbf{Вход:} Модель $M_{net}$, Координаты $C \in \mathbb{R}^{1 \times N \times 3 \times 3}$, Маска $M \in \mathbb{B}^{1 \times N}$, Число проходов $M_{runs}$
\Statex \textbf{Выход:} Средняя вероятность свертываемости $\bar{p}$, Эпистемическая неопределенность $U$
\State Set model to evaluation mode: $M_{net}.\operatorname{eval}()$
\State Compute heavy biophysical features (frontend):
\State \quad $Feats \gets M_{net}.\operatorname{compute\_features}(C, M)$ \Comment{выполняется строго один раз}
\State Save training flags of all network modules
\State \textbf{Enable Dropout during inference:}
\For{each module $m \in M_{net}.\operatorname{modules}()$}
    \If{$m$ is an instance of \texttt{nn.Dropout}}
        \State Set module to training mode: $m.\operatorname{train}()$
    \EndIf
\EndFor
\State Initialize list of folding probabilities: $Probs \gets []$
\For{$run \gets 1$ \textbf{to} $M_{runs}$}
    \State Run encoder forward pass: $H_{node}, \_ \gets M_{net}.\operatorname{encode\_features}(Feats)$
    \State Run heads prediction: $P \gets M_{net}.\operatorname{predict\_from\_repr}(H_{node})$
    \State Append to list: $Probs.\operatorname{append}(\sigma(P_{fold\_logit}))$
\EndFor
\State Restore saved training flags of all modules
\State Compute mean and variance statistics over runs:
\State \quad $\bar{p} \gets \operatorname{mean}(Probs)$
\State \quad $U \gets \operatorname{variance}(Probs)$
\State \Return $\bar{p}, U$
\end{algorithmic}
\end{algorithm}

\newpage

% --- ПОДГОТОВКА ДАННЫХ И ЭКСПЕРИМЕНТАЛЬНОЕ ИССЛЕДОВАНИЕ ---
\section{Подготовка данных и экспериментальное исследование}

\subsection{Формирование набора данных и предотвращение утечки информации}
Для обучения и объективной валидации модели был собран сбалансированный набор данных, состоящий из нативных структур и синтетических структур-декоев с различными уровнями геометрических искажений.
Нативные структуры (положительный класс, \texttt{label} $= 1.0$) были получены из базы данных PDB (Protein Data Bank) с помощью специализированного запроса к RCSB API. Фильтрация проводилась по следующим жестким биофизическим критериям:
\begin{itemize}
    \item Метод определения пространственной структуры: рентгеноструктурный анализ (X-ray diffraction).
    \item Разрешение структуры: $\le 2.0\,\text{\AA}$.
    \item Длина белковой цепи: от $50$ до $700$ аминокислотных остатков.
    \item Непрерывность цепи: максимальное расстояние между соседними атомами остова $C_i - N_{i+1}$ не превышает $2.0\,\text{\AA}$.
\end{itemize}
В результате было отобрано $141\,726$ стабильных белковых цепей. Для каждой цепи было сгенерировано по два декоя, что в сумме составило $425\,178$ образцов ($297\,624$ / $63\,774$ / $63\,780$ в обучающей, валидационной и тестовой подвыборках соответственно).

Для предотвращения утечки информации (data leakage), представляющей серьезную методологическую проблему при машинном обучении в структурной биологии, разбиение на подвыборки (обучающую, валидационную и тестовую в соотношении 70\% / 15\% / 15\%) проводилось строго по PDB ID и Chain ID (идентификатору родительской цепи белка). Это гарантирует, что исходная нативная структура цепи белка и все сгенерированные на ее основе декои всегда попадают в один и тот же сплит. При случайном поэлементном разбиении схожие конформации одной и той же цепи белка распределились бы между обучением и тестом, что привело бы к искусственному завышению метрик точности модели.

\subsection{Стратегии негативного сэмплирования (генерация декоев)}
Для каждого нативного белка было сгенерировано два декой-структурных ловушки с использованием семи различных деформационных стратегий, имитирующих ошибки современных генераторов белковых остовов (RFDiffusion, Chroma и др.):
\begin{enumerate}
    \item \textbf{easy\_global\_noise} (простой декой, класс 1): наложение независимого гауссова шума со среднеквадратичным отклонением $1.5\,\text{\AA}$ на координаты всех атомов остова, что соответствует общему температурному или физическому коллапсу цепи.
    \item \textbf{easy\_chain\_break} (простой декой, класс 1): разрез остова в случайном месте с трансляционным смещением одной из половин цепи на расстояние от $12$ до $18\,\text{\AA}$, что позволяет модели детектировать разрывы пептидной связи.
    \item \textbf{hard\_core\_unpacked} (сложный декой, класс 2): радиальное расширение (раздутие) цепи от центра масс на 10--25\%, что нарушает плотную упаковку гидрофобного ядра.
    \item \textbf{hard\_false\_compact} (сложный декой, класс 2): разделение остова на 2--4 сегмента с поворотом каждого блока на случайный угол от $20^\circ$ до $120^\circ$ вокруг своей оси и последующей плотной упаковкой в границах сферы малого радиуса (4.5--8.5\,\text{\AA}). Такое искажение имитирует ложно-компактные структуры без упорядоченной геометрии.
    \item \textbf{hard\_near\_native} (сложный декой, класс 3): анизотропное масштабирование цепи вдоль координатных осей (сжатие/растяжение на 2--5\%) с последующими микро-ротациями, что создает структуры, близкие к нативным по расстояниям, но термодинамически нестабильные.
    \item \textbf{borderline\_hinge\_defect} (пограничный декой, класс 4): вращение концевого фрагмента цепи вокруг случайной поворотной точки (шарнира) на угол от $8^\circ$ до $22^\circ$, что моделирует шарнирный дефект сборки доменов.
    \item \textbf{borderline\_local\_fragment\_rotation} (пограничный декой, класс 4): выборка короткого фрагмента (3--8 остатков) и его поворот на угол $6^\circ$--$18^\circ$ вокруг вектора, соединяющего концы фрагмента. Имитирует локальные нарушения торсионных углов в петлях.
\end{enumerate}

Для всех образцов были автономно вычислены точные физические таргеты. Стерические конфликты детектировались по виртуальным атомам $C_\beta$ (порог $3.5\,\text{\AA}$); пары остатков с $|i-j| < 3$ исключались из подсчета, поскольку их близость обусловлена геометрией самой цепи, а не стерическим конфликтом. Таргеты водородных связей в финальной версии модели не используются (см. разд. 1.2).

\subsection{Гиперпараметры модели и параметры обучения}
В табл. \ref{tab:hyperparams} представлены гиперпараметры модели, используемые в ходе экспериментального исследования.

\begin{table}[h]
\centering
\caption{Гиперпараметры нейросетевой модели BbValidator}
\label{tab:hyperparams}
\begin{tabular}{lc}
\toprule
\textbf{Параметр} & \textbf{Значение} \\
\midrule
Размерность скрытого представления ($d_{model}$) & 192 \\
Количество MPNN-слоев & 2 \\
Количество слоев Transformer Encoder & 4 \\
Количество голов внимания в трансформере & 8 \\
Размерность скрытого слоя ребер ($pair\_dim$) & 64 \\
Вероятность регуляризации Dropout (энкодер и головы) & 0.15 \\
Размер батча (Batch Size) & 64 \\
Количество эпох обучения & 15 \\
Скорость обучения AdamW (веса модели) & $6 \times 10^{-4}$ \\
Скорость обучения AdamW (параметры лосса $s_t$) & $10^{-3}$ \\
Коэффициент затухания весов (Weight Decay) & $10^{-4}$ \\
Глобальный сид воспроизводимости (torch/numpy/сэмплер) & 42 \\
Бакетизация батчей по длине цепи (шаг) & 50 остатков \\
Веса классов failure mode & обратная частота в train \\
Критерий выбора лучшего чекпоинта & $(\mathrm{AUC} + \mathrm{PR\text{-}AUC})/2$ на val \\
\bottomrule
\end{tabular}
\end{table}

Обучение проводилось с автоматическим смешиванием точности (AMP в формате \texttt{bfloat16}) для ускорения вычислений на GPU и экономии памяти. В графовых слоях MPNN использовался механизм градиентного чекпоинтинга (\texttt{gradient}\-\texttt{checkpointing}), что позволило обучать модель на длинных белковых цепях без переполнения памяти ускорителя.

\subsection{Анализ сходимости и результаты обучения}
Обучение проводилось в течение 15 эпох с выбором лучшего чекпоинта по составной метрике $(\mathrm{AUC} + \mathrm{PR\text{-}AUC})/2$ на валидационной подвыборке; все приведенные ниже результаты получены на тестовой подвыборке ($63\,780$ образцов), которая не использовалась ни при обучении, ни при выборе чекпоинта. Модель показала монотонную сходимость: точность на обучении достигла $95.3\%$, на валидации --- $94.7\%$ (AUC $0.979$) к последней эпохе без признаков переобучения.

Для оценки вклада архитектуры энкодера были обучены две базлайновые модели по идентичному протоколу (тот же биофизический фронтенд, пулинг, головы, функция потерь, сиды и критерий селекции; различается только энкодер, $d_{model}=128$):
\begin{itemize}
    \item \textbf{MLP} --- независимый многослойный перцептрон на каждый остаток, без какого-либо обмена информацией между позициями;
    \item \textbf{GPS} --- стек слоев GPSConv (графовое внимание TransformerConv + глобальное внимание) на том же kNN-графе фронтенда.
\end{itemize}

\begin{table}[h]
\centering
\caption{Результаты на тестовой подвыборке: главная задача (foldability)}
\label{tab:test_results}
\begin{tabular}{lccc}
\toprule
\textbf{Метрика} & \textbf{Гибрид (d=192)} & \textbf{GPS (d=128)} & \textbf{MLP (d=128)} \\
\midrule
Accuracy & 0.9470 & 0.9461 & 0.9001 \\
ROC-AUC & 0.9789 & \textbf{0.9798} & 0.9497 \\
PR-AUC & 0.9357 & \textbf{0.9408} & 0.8545 \\
ECE (калибровка) & 0.0098 & \textbf{0.0079} & 0.0228 \\
Precision & 0.8796 & \textbf{0.8858} & 0.7840 \\
Recall & \textbf{0.9743} & 0.9625 & 0.9666 \\
Failure mode accuracy & \textbf{0.9426} & 0.9401 & 0.8950 \\
\bottomrule
\end{tabular}
\end{table}

Обе модели с обменом информацией между остатками (гибрид и GPS) существенно превосходят MLP по всем метрикам, что подтверждает важность пространственного контекста. При этом статистически значимого преимущества гибридной архитектуры над чистым GPS на данном наборе данных не выявлено: различия ранговых метрик лежат в пределах $0.1$ п.п. Гибрид демонстрирует несколько лучший recall нативных структур и точность классификации типа дефекта. Все модели хорошо калиброваны (ECE $< 0.023$), что обосновывает использование вероятностных предсказаний и MC-Dropout неопределенности на практике. Тривиальный предсказатель <<всегда декоев>> (с учетом дисбаланса классов 1:2) дает точность $66.7\%$.

\begin{table}[h]
\centering
\caption{Специфичность по стратегиям декоев (доля корректно отвергнутых) и AUC <<нативы против семейства>> на тестовой подвыборке}
\label{tab:per_strategy}
\begin{tabular}{lcccccc}
\toprule
 & \multicolumn{2}{c}{\textbf{Гибрид}} & \multicolumn{2}{c}{\textbf{GPS}} & \multicolumn{2}{c}{\textbf{MLP}} \\
\cmidrule(lr){2-3} \cmidrule(lr){4-5} \cmidrule(lr){6-7}
\textbf{Стратегия} & spec & AUC & spec & AUC & spec & AUC \\
\midrule
easy\_global\_noise & 1.000 & 1.000 & 1.000 & 1.000 & 1.000 & 1.000 \\
easy\_chain\_break & 1.000 & 1.000 & 1.000 & 1.000 & 1.000 & 1.000 \\
hard\_core\_unpacked & 1.000 & 1.000 & 1.000 & 1.000 & 0.998 & 1.000 \\
hard\_false\_compact & 1.000 & 1.000 & 1.000 & 1.000 & 1.000 & 1.000 \\
hard\_near\_native & 0.999 & 1.000 & 1.000 & 1.000 & 0.995 & 0.999 \\
borderline\_hinge\_defect & 0.830 & 0.947 & 0.847 & 0.951 & 0.716 & 0.909 \\
borderline\_local\_fragment\_rotation & 0.726 & 0.912 & 0.739 & 0.914 & 0.400 & 0.756 \\
\bottomrule
\end{tabular}
\end{table}

Разбивка по стратегиям (табл. \ref{tab:per_strategy}) показывает, что все категории <<простых>> и <<сложных>> декоев детектируются практически безошибочно, тогда как узким местом являются пограничные локальные дефекты --- в первую очередь повороты коротких фрагментов, изменяющие геометрию лишь в небольшой области цепи. Это согласуется с физической интуицией: локальный дефект слабо влияет на глобальные дескрипторы структуры.

\subsection{Валидация на выходах внешних генераторов (OOD)}
Для проверки обобщающей способности за пределы обучающего распределения (все декои обучения --- синтетические деформации нативных структур) модели были оценены на $18\,000$ реально сгенерированных структур из открытого бенчмарка MotifBench \cite{motifbench} (выходы пяти генераторов: RFdiffusion, RFdiffusion-AA, ODesign-Rigid, GPDL, EvoDiff), не встречавшихся ни в обучении, ни среди декоев.

\begin{table}[h]
\centering
\caption{Средний $P(\mathrm{Foldable})$ на выходах генераторов (MotifBench, $n=3000$ на генератор). Для сравнения: нативы тестовой выборки --- $0.87/0.76/0.73$, синтетические декои --- $0.10/0.08/0.16$ для гибрида/GPS/MLP}
\label{tab:ood}
\begin{tabular}{lcccc}
\toprule
\textbf{Генератор} & \textbf{scRMSD, \AA} & \textbf{Гибрид} & \textbf{GPS} & \textbf{MLP} \\
\midrule
RFdiffusion & 4.70 & \textbf{0.845} & 0.658 & 0.704 \\
ODesign-Rigid & 5.14 & \textbf{0.761} & 0.685 & 0.661 \\
RFdiffusion-AA & 6.25 & \textbf{0.689} & 0.564 & 0.647 \\
GPDL & 6.35 & \textbf{0.532} & 0.450 & 0.273 \\
EvoDiff & 12.31 & 0.384 & 0.401 & \textbf{0.084} \\
\bottomrule
\end{tabular}
\end{table}

Все три архитектуры одинаково ранжируют генераторы по качеству, согласуясь с независимой ground-truth метрикой MotifBench --- само-согласованным scRMSD (средний RMSD ESMFold-рефолдов спроектированных последовательностей на сгенерированный остов, табл. \ref{tab:ood}). Корреляция Спирмена $P(\mathrm{Foldable}) \leftrightarrow \mathrm{scRMSD}$ значима ($p \le 10^{-11}$) и отрицательна на четырех генераторах из пяти для всех моделей: от $-0.17$ до $-0.29$ у гибрида, до $-0.46$ у MLP на GPDL. Лишь на выходах EvoDiff (худший генератор, $3.1\%$ образцов с scRMSD $< 2\,\text{\AA}$) поведение моделей расходится: MLP сохраняет правильный знак корреляции ($-0.33$), гибрид становится неразличимым, а GPS инвертирует его. Таким образом, межостаточные взаимодействия дают большой выигрыш внутри обучающего распределения, однако простой локальный скоринг (MLP) оказывается наиболее робастным на дальнем OOD --- классический компромисс емкости и устойчивости к сдвигу распределения.

Отдельный кейс: на $29$ структурах, сгенерированных ранним (пилотным) чекпоинтом SE(3)-диффузионной модели motif-scaffolding, все три модели единогласно дали $P(\mathrm{Foldable}) \approx 0$. Физическая диагностика подтвердила корректность вердикта: образцы содержали $27\%$ выбросов карты Рамачандрана и $0.49$ стерических конфликта на остаток при норме $1$--$2\%$ и $\sim 0.03$ соответственно. Данный результат демонстрирует применимость модели как фильтра физически некорректных генераций за пределами обучающего распределения.

\subsection{Оценка производительности (Бенчмарк скорости)}
Оценка вычислительной производительности модели проводилась с помощью скрипта \texttt{benchmark.py}. Алгоритм тестирования включает прогревочный цикл из 5 итераций, после чего измеряются задержки инференса на протяжении 50 проходов. Измерения проводятся для различного количества проходов MC-Dropout ($M \in \{8, 16, 32\}$).

Поскольку тяжелые физические дескрипторы (фронтенд) вычисляются строго один раз для каждого белка и не требуют повторного прохождения для каждого MC-прохода, инференс модели является высокопроизводительным. Модель способна обрабатывать от десятков до сотен белковых цепей в секунду (в зависимости от длины цепи и режима работы с MC-Dropout), что делает ее применимой для фильтрации миллионов сгенерированных белковых остовов de novo в реальном времени.

\newpage

% --- ЗАКЛЮЧЕНИЕ ---
\section{Заключение}
В ходе выполнения данной курсовой работы была спроектирована, реализована и исследована нейросетевая модель \texttt{BbValidator} для валидации белковых структур остова белка. 

Основные результаты работы:
\begin{enumerate}
    \item Разработан эффективный физический фронтенд \texttt{Biophysical}\allowbreak\texttt{Frontend}, позволяющий вычислять 31 признак узлов и 20 признаков пар на основе фи\-зи\-ко-ге\-о\-мет\-ри\-чес\-ких законов укладки белковых цепей, включая PCA-проекцию фрагментных расстояний, обучаемую по 31.5 млн фрагментов нативных структур.
    \item Предложена и реализована гибридная архитектура \texttt{Hybrid}\allowbreak\texttt{Protein}\allowbreak\texttt{Encoder}, сочетающая преимущества графовых сетей (MPNN) в обработке локального окружения и трансформеров в анализе глобального контекста.
    \item Внедрен механизм автоматического балансирования потерь на основе гомоскедастической неопределенности, что позволило стабилизировать обучение без ручного подбора весов задач.
    \item Реализован метод Монте-Карло Дропаута (\texttt{MC-Dropout}) на этапе инференса для количественной оценки эпистемической неопределенности предсказаний.
    \item Проведено экспериментальное исследование на наборе из $425\,178$ образцов: модель достигла ROC-AUC $0.979$ и точности $94.7\%$ на тестовой подвыборке при калибровке ECE $0.0098$; сравнение с базлайнами (MLP, GPSConv) показало, что обмен информацией между остатками дает основной прирост качества, при этом пограничные локальные дефекты остаются наиболее трудным классом.
    \item Выполнена валидация на выходах пяти внешних генераторов белковых остовов ($18\,000$ структур, бенчмарк MotifBench): предсказания модели значимо коррелируют с независимой метрикой само-согласованности scRMSD и корректно ранжируют генераторы по качеству, включая структуры, принципиально отсутствовавшие в обучающем распределении.
\end{enumerate}

Разработанный конвейер готов к интеграции в современные генеративные пайплайны дизайна белков de novo в качестве модуля ранней фильтрации сгенерированных пептидных остовов.

\newpage

% --- СПИСОК ЛИТЕРАТУРЫ ---
\begin{thebibliography}{99}
\addcontentsline{toc}{section}{Список литературы}

\bibitem{alphafold} 
Jumper, J., Evans, R., Pritzel, A. et al. Highly accurate protein structure prediction with AlphaFold. \textit{Nature} 596, 583–589 (2021).

\bibitem{rfdiffusion} 
Watson, J. L., Juergens, D., Bennett, N. R. et al. De novo design of protein structure and function with RFdiffusion. \textit{Nature} 620, 1089–1100 (2023).

\bibitem{chroma} 
Ingraham, J., Barzilay, R., Regina, A. et al. Illuminating protein space with a programmable generative model. \textit{bioRxiv} 2022.12.01.518682 (2022).

\bibitem{kendall2018} 
Kendall, A., Gal, Y., Cipolla, R. Multi-task learning using uncertainty to weigh losses for scene geometry and semantics. In \textit{Proceedings of the IEEE conference on computer vision and pattern recognition}, 7482–7491 (2018).

\bibitem{gal2016} 
Gal, Y., Ghahramani, Z. Dropout as a bayesian approximation: Representing model uncertainty in deep learning. In \textit{International conference on machine learning}, 1050–1059 (2016).

\bibitem{pyg} 
Fey, M., Lenssen, J. E. Fast graph representation learning with PyTorch Geometric. \textit{arXiv preprint arXiv:1903.02428} (2019).

\bibitem{biotite}
Kunzmann, P., Hamacher, K. Biotite: a unifying framework for biophysical data in Python. \textit{BMC Bioinformatics} 19, 1–10 (2018).

\bibitem{motifbench}
Alamdari, S., Yang, K. MotifBench: результаты оценки генераторов белковых остовов (motif scaffolding) [датасет]. \textit{Zenodo} (2025). DOI: 10.5281/zenodo.15445142 и связанные записи серии.

\end{thebibliography}

\end{document}
